% ============================================================
% Capitolo 14 — Multi-agente
% ============================================================

\chapter{Multi-agente: Claude, Copilot, Codex, Gemini}
\label{ch:multi-agente}

AI-DLC è progettato per funzionare con più agenti, simultaneamente o in sequenza, sullo stesso progetto, con lo stesso \texttt{\_CONTEXT.md} come fonte di verità condivisa.

\section{Uno stesso progetto, più entry point}
\label{sec:14-entry-point}

\begin{table}[htbp]
  \centering
  \begin{tabular}{lll}
    \toprule
    \textbf{Agente} & \textbf{File di startup} & \textbf{Note} \\
    \midrule
    Claude Code & \texttt{CLAUDE.md} & Importa \texttt{AGENTS.md} \\
    OpenAI Codex CLI & \texttt{AGENTS.md} & Entry point principale \\
    Google Gemini & \texttt{GEMINI.md} & Delega ad \texttt{AGENTS.md} \\
    OpenClaw & \texttt{OPENCLAW.md} & Delega ad \texttt{AGENTS.md} \\
    GitHub Copilot & \texttt{.github/copilot-instructions.md} & Non può importare \\
    Visual Studio 2026 & \texttt{VISUALSTUDIO.md} & Importa \texttt{AGENTS.md}; regola DOC-01 \\
    \bottomrule
  \end{tabular}
  \caption{Entry point per ogni agente AI}
  \label{tab:14-entry-points}
\end{table}

La struttura è a stella: \texttt{AGENTS.md} è il centro, gli altri file orbitano intorno. Il contenuto sostanziale del protocollo --- le regole di rischio, i confidence tag, gli HALT trigger --- vive in \texttt{AGENTS.md}.

\diagramfig{ch14-stella-multiagente}{Architettura multi-agente a stella: \texttt{AGENTS.md} al centro, \texttt{\_CONTEXT.md} come fonte di verità condivisa.}

\begin{tipbox}
\textbf{Regola d'oro:} modifica solo \texttt{AGENTS.md} per le regole comuni. Non modificare \texttt{CLAUDE.md}, \texttt{GEMINI.md}, \texttt{OPENCLAW.md} per aggiungere regole --- aggiornali solo per comportamenti specifici di quell'agente.
\end{tipbox}

\section{Un workflow multi-agente su TaskFlow API}
\label{sec:14-workflow}

\begin{itemize}
  \item \textbf{Lorenzo} usa Claude Code per design e architettura
  \item \textbf{Giulia} usa GitHub Copilot in VS Code per lo sviluppo quotidiano
  \item \textbf{Marco} usa Claude Code per code review, Codex CLI per script
\end{itemize}

Il \texttt{\_CONTEXT.md} è condiviso in git. Ogni mattina, quando Lorenzo fa il bootstrap, Claude Code legge lo stesso file che Giulia ha aggiornato la sera prima. Il contesto è allineato senza sincronizzazione manuale.

\section{Mantenere Copilot allineato}
\label{sec:14-sync-copilot}

Copilot non può importare file esterni, quindi il suo file di istruzioni deve essere mantenuto allineato con \texttt{AGENTS.md}:

\begin{lstlisting}[language=bash,caption={Verifica allineamento Copilot}]
bash .ai-dlc/tools/sync-copilot.sh
\end{lstlisting}

\begin{lstlisting}[language=,caption={Output di sync-copilot}]
[OK] Risk classification: present in both
[OK] HALT trigger reference: present in both
[!]  @alternatives command: present in AGENTS.md,
   missing in copilot-instructions.md
\end{lstlisting}

Esegui \texttt{sync-copilot} dopo ogni modifica a \texttt{AGENTS.md} e prima di aggiungere un nuovo membro al team che usa Copilot.

\section{I limiti del multi-agente}
\label{sec:14-limiti}

\textbf{Non è collaborazione in tempo reale.} Gli agenti non si parlano tra loro. La coerenza è garantita dal \texttt{\_CONTEXT.md} condiviso, non da un canale diretto. Conflitti git su edit concorrenti rimangono conflitti git.

\textbf{La qualità del context condiviso determina la qualità del risultato.} Se \texttt{\_CONTEXT.md} è obsoleto, tutti gli agenti lavorano su informazioni sbagliate.

\section{Visual Studio 2026 e la regola DOC-01}
\label{sec:14-visualstudio}

La versione 4.0.0 introduce \texttt{VISUALSTUDIO.md}, il file di startup per Visual Studio 2026. Importa \texttt{AGENTS.md} e aggiunge le specificità dell'ambiente solution/project (\texttt{.sln} e \texttt{.csproj}).

La regola \textbf{DOC-01} impone: ogni progetto ha la propria documentazione in \texttt{projects\textbackslash <ProjectName>\textbackslash} (nella radice del repository). Nessun file di documentazione va nelle cartelle del codice sorgente.

La solution include un progetto condiviso \texttt{AI-DLC.Documentation}: il file \texttt{AI-DLC.Documentation.csproj} sta nella radice del repository e raccoglie le cartelle di documentazione, anch'esse in radice:
\begin{itemize}
  \item \texttt{projects/<Project>/\_CONTEXT.md} --- stato del progetto
  \item \texttt{projects/<Project>/PROGRESS.md} --- progress log
  \item \texttt{projects/<Project>/instructions.md} --- regole specifiche
  \item \texttt{projects/shared/} --- documentazione cross-progetto
  \item \texttt{docs/\_solution/} --- documentazione a livello di solution
\end{itemize}

Non è necessario usare \texttt{VISUALSTUDIO.md} se il progetto non ha una solution Visual Studio --- è completamente opzionale.

\section*{\LabelSummary}

\begin{itemize}
  \item \texttt{AGENTS.md} è il centro del sistema multi-agente; gli altri file sono wrapper.
  \item Il \texttt{\_CONTEXT.md} condiviso in git garantisce la coerenza tra agenti diversi.
  \item \texttt{sync-copilot} verifica l'allineamento di Copilot dopo ogni modifica alle regole comuni.
  \item Visual Studio 2026 introduce la regola DOC-01: documentazione separata dal codice in \texttt{AI-DLC.Documentation}.
  \item Il multi-agente funziona bene per lavoro asincrono; non risolve i conflitti di edit concorrenti.
\end{itemize}
